\documentclass[11pt]{article}

\usepackage[T1]{fontenc}
\usepackage[utf8]{inputenc}
\usepackage{lmodern}
\usepackage[margin=2.5cm]{geometry}
\usepackage{amsmath,amssymb}
\usepackage{booktabs}
\usepackage{tabularx}
\usepackage{array}
\usepackage{enumitem}
\usepackage{xcolor}
\usepackage{fancyvrb}
\usepackage[hidelinks]{hyperref}

\newcolumntype{L}{>{\raggedright\arraybackslash}X}
\renewcommand{\arraystretch}{1.15}
\setlist{itemsep=2pt,topsep=4pt}
\newcommand{\var}[1]{\texttt{#1}}

\title{Entropy Meets Topology:\\[2pt]
\large Transferring Atrial Fibrillation Detection from Clinical ECG to Wrist-Worn PPG\\
via Nonlinear Dynamical Features}
\author{MATH432 Mathematical Data Science --- Final Project}
\date{Draft: Sections 1--2}

\begin{document}
\maketitle

%=====================================================================
\section{Introduction}
%=====================================================================

\subsection{Background}

Atrial fibrillation (AF) is the most common sustained cardiac arrhythmia.
During AF, the atria do not contract in an organised way. Electrical impulses
reach the ventricles at irregular times, so \emph{the interval between
consecutive heartbeats becomes irregular}. AF is associated with an
approximately five-fold increase in the risk of stroke~\cite{wolf1991}. It is
also often paroxysmal (it comes and goes) and asymptomatic, so a short ECG
recorded in a clinic can easily miss it. For these reasons, long-term and
low-cost rhythm monitoring is clinically valuable.

Consumer wearables now make such monitoring possible. A smartwatch carries a
\emph{photoplethysmography} (PPG) sensor: an LED measures changes in blood
volume under the skin of the wrist. From this signal the watch can estimate
the time between pulses passively and continuously. Large studies such as the
Apple Heart Study~\cite{perez2019} have shown that PPG-based irregular-rhythm
notifications can be used for AF screening at population scale. However, two
problems make AF detection from wrist PPG harder than from ECG:
\begin{enumerate}
  \item \textbf{Scarce labels.} Annotating AF requires a simultaneous
  reference ECG reviewed by clinicians. There are therefore far fewer
  labelled wearable recordings than labelled clinical ECG recordings, such as
  the MIT-BIH databases on PhysioNet~\cite{physionet}.
  \item \textbf{Domain shift.} A PPG pulse is smoother than an ECG QRS
  complex, so beat timing from PPG is less precise. PPG is also very
  sensitive to motion artifacts. A model trained on clinical ECG may
  therefore fail on wrist PPG.
\end{enumerate}

A natural bridge between the two domains is the \emph{beat-to-beat interval
series}. From ECG we obtain RR intervals (between R-peaks); from PPG we obtain
pulse-to-pulse intervals (PPI). Both measure essentially the same quantity,
the time from one heartbeat to the next. Since the defining symptom of AF is
the irregularity of exactly this series, features computed from it should
keep their meaning across domains.

The interval series is the output of a nonlinear dynamical system (the cardiac
conduction system under autonomic control). Classical HRV statistics, such as
the standard deviation or RMSSD, summarise only its spread. Three families of
mathematical tools describe its \emph{dynamics} instead:
\begin{itemize}
  \item \textbf{Scaling laws.} Detrended fluctuation analysis
  (DFA)~\cite{peng1995} measures the fractal correlation of the series.
  Healthy sinus rhythm shows $1/f$-type scaling (exponent $\alpha \approx 1$),
  while an uncorrelated series gives $\alpha \approx 0.5$.
  \item \textbf{Entropy.} Sample entropy~\cite{richman2000} and its
  AF-specific variant COSEn~\cite{lake2011} quantify how unpredictable the
  next interval is. Permutation entropy~\cite{bandt2002}, computed at several
  time scales~\cite{costa2002}, measures the diversity of up/down ordering
  patterns.
  \item \textbf{Topology.} Takens' theorem~\cite{takens1981} says that a
  delay embedding of a scalar series reconstructs the geometry of the
  underlying attractor. Persistent homology~\cite{edelsbrunner2010} then
  describes the shape of the resulting point cloud (connected components,
  loops) at all scales at once. The stability theorem~\cite{cohensteiner2007}
  guarantees that this description changes only a little when the data are
  perturbed a little, which is an attractive property for noisy wearable
  signals. Persistence images~\cite{adams2017} and persistence
  landscapes~\cite{bubenik2015} turn these descriptions into vectors that a
  classifier can use.
\end{itemize}

\subsection{Problem statement}

This project asks whether \textbf{AF detection learned from clinical ECG can
be transferred to wrist-worn PPG} by using mathematically motivated features
of the beat-interval series. We formulate four research questions
(Table~\ref{tab:rq}).

\begin{table}[ht]
\centering
\caption{Research questions.}
\label{tab:rq}
\begin{tabularx}{\textwidth}{@{}lL@{}}
\toprule
 & Research question \\
\midrule
\textbf{RQ1} & How well do scaling, entropy and topological features
distinguish AF from non-AF on clinical ECG? Are the three families
complementary to each other, and do they add information beyond standard HRV
statistics? \\
\textbf{RQ2} & Does a model trained on clinical ECG still work on wrist PPG?
How large is the domain gap? \\
\textbf{RQ3} & When only a few patients with labelled wrist data are
available, does adding clinical ECG data improve performance? \\
\textbf{RQ4} & Which features remain reliable under real motion artifacts
recorded during daily life? \\
\bottomrule
\end{tabularx}
\end{table}

The task is binary classification of fixed-length segments of the interval
series into \emph{AF} and \emph{non-AF}. All evaluations are
patient-independent: no patient contributes data to both the training and the
test set.

\subsection{Why this problem matters}

\begin{itemize}
  \item \textbf{Clinical relevance.} Earlier detection of asymptomatic AF can
  lead to earlier treatment and prevention of stroke. Wrist PPG is the only
  signal that a wearable can record continuously without any action from the
  user.
  \item \textbf{Data efficiency.} If clinical ECG databases, which are public
  and well annotated, can compensate for scarce wearable labels, new wearable
  devices could be deployed with much less data collection.
  \item \textbf{Interpretability and robustness.} Commercial detection
  algorithms are proprietary, and recent deep-learning approaches are hard to
  interpret. Each feature in this project has a clear mathematical meaning,
  for example ``the rhythm has lost its fractal correlation'' or ``the next
  interval is unpredictable''. The stability theorem of persistent homology
  also gives a theoretical reason to expect robustness to noise, which we
  test directly in RQ4.
  \item \textbf{Connection to the course.} The project combines empirical
  scaling laws, entropy theory, topological data analysis and classification
  on one real-world problem.
\end{itemize}

%=====================================================================
\section{Data}
%=====================================================================

We use two public datasets (Table~\ref{tab:datasets}). Both are reduced to
the same representation: a series of beat intervals in seconds, with an AF /
non-AF label for each interval.

\begin{table}[ht]
\centering
\caption{Datasets used in this project.}
\label{tab:datasets}
\begin{tabularx}{\textwidth}{@{}LLLL@{}}
\toprule
Dataset & Signal & Size & Role \\
\midrule
MIT-BIH Atrial Fibrillation Database (AFDB)~\cite{physionet,moody1983}
  & 2-lead Holter ECG, 250\,Hz
  & 25 patients, about 10\,h each
  & Stage~1: feature comparison on clinical ECG (RQ1); source domain for
    transfer (RQ2, RQ3) \\
\addlinespace
Long-term ECG and wrist PPG with annotated AF episodes (Zenodo,
v6)~\cite{bacevicius2026}
  & Wrist PPG 100\,Hz, reference ECG 500\,Hz, 3-axis accelerometers
  & 39 patients (35 with PPG), 5--8 days each, 261 days in total
  & Stages~2--3: target domain for transfer (RQ2, RQ3); motion robustness
    (RQ4) \\
\bottomrule
\end{tabularx}
\end{table}

%---------------------------------------------------------------------
\subsection{MIT-BIH Atrial Fibrillation Database (AFDB)}
%---------------------------------------------------------------------

\paragraph{Source.} PhysioNet,
\url{https://physionet.org/content/afdb/1.0.0/} (open access). We download it
from PhysioNet's official open-data mirror and verify every file against the
published SHA-256 checksums.

\paragraph{Content.} The database contains 25 long-term ECG recordings from
patients with AF, mostly paroxysmal. Each recording is about 10 hours long and
has two ECG leads sampled at 250\,Hz with 12-bit resolution. Two recordings
(00735, 03665) contain annotations only, with no signal. This does not affect
us, because we use only the beat times.

\paragraph{Raw file structure.} WFDB format, one set of files per record
(Table~\ref{tab:afdb-files}).

\begin{table}[ht]
\centering
\caption{AFDB files for one record.}
\label{tab:afdb-files}
\begin{tabularx}{\textwidth}{@{}lL@{}}
\toprule
File & Content \\
\midrule
\var{.hea} & Header: number of channels, sampling frequency, number of
samples, gain \\
\var{.dat} & ECG signal (binary, two channels) \\
\var{.atr} & Rhythm annotations made by cardiologists: the sample index where
each rhythm episode begins, with a label \var{(N} (normal), \var{(AFIB} (AF),
\var{(AFL} (atrial flutter) or \var{(J} (AV junctional rhythm) \\
\var{.qrs} & Beat (QRS) locations from an automatic detector \\
\var{.qrsc} & Manually corrected beat locations (only records 05091 and
07859) \\
\bottomrule
\end{tabularx}
\end{table}

\noindent\textit{Sample --- record 04746:}
\begin{Verbatim}[fontsize=\small,frame=single,framesep=6pt]
04746.hea   04746 2 250 9205760 22:00:00     # 2 channels, 250 Hz, 10.2 h
            04746.dat 212 0 12 0 -57 10750 0 ECG1
            04746.dat 212 0 12 0 -68 -18834 0 ECG2

04746.dat   (first 5 samples, mV)   ECG1    ECG2
                                   -0.285  -0.340
                                   -0.300  -0.320
                                   -0.310  -0.325
                                   -0.300  -0.315
                                   -0.300  -0.310

04746.atr   sample      rhythm               # rhythm changes
            62          (N
            1,097,510   (AFIB
            1,098,684   (N
            2,284,261   (AFIB

04746.qrs   sample: 124, 357, 583, 809, 1033, 1263, ...   # beat locations
\end{Verbatim}

\paragraph{From annotations to labelled intervals.} For each record:
\begin{enumerate}
  \item take the beat locations from \var{.qrsc} if available, otherwise from
  \var{.qrs};
  \item compute RR intervals,
  $\mathrm{rr}_i = (b_{i+1} - b_i)/250$ seconds, where $b_i$ is the sample
  index of beat $i$;
  \item give each interval the rhythm label in force at the beat that ends it.
\end{enumerate}
Intervals outside $[0.25, 2.5]$\,s are flagged as implausible. In total this
gives 1{,}221{,}650 intervals (249\,h). By time, 59.9\% is normal rhythm,
37.5\% AF, 2.1\% junctional rhythm and 0.6\% atrial flutter.

\paragraph{Segment dataset.} The interval series is cut into non-overlapping
segments of \textbf{60 consecutive intervals}. We use a fixed number of beats
rather than a fixed duration, so that every segment gives the same number of
points to the entropy and topology computations regardless of heart rate. The
labelling rules are:
\begin{itemize}
  \item A segment is dropped if it contains any implausible interval, any
  atrial flutter, or any interval without a rhythm label.
  \item AF is the AFIB label only. Normal and junctional rhythm count as
  non-AF.
  \item Main label: $\var{label\_50} = 1$ if at least 50\% of the segment's
  intervals are AF.
  \item Secondary label for a sensitivity analysis: $\var{label\_excl} = 1$
  if at least 90\% AF, $0$ if at most 10\% AF, otherwise excluded ($-1$).
\end{itemize}
Record 07162 is excluded. Its reference beat locations do not match the QRS
complexes in its signal: an independent detector agreed with only about 30\%
of its beats, while the other records agreed at around 99\%.

The resulting dataset has \textbf{18{,}925 segments from 24 patients}. Of
these, 7{,}724 (40.8\%) are AF under \var{label\_50}; 387 segments (2.0\%)
are ambiguous under \var{label\_excl}. Its fields are listed in
Table~\ref{tab:afdb-fields}, and two example segments in
Table~\ref{tab:afdb-sample}.

\begin{table}[ht]
\centering
\caption{Fields of the AFDB segment dataset.}
\label{tab:afdb-fields}
\begin{tabularx}{\textwidth}{@{}llL@{}}
\toprule
Field & Shape & Meaning \\
\midrule
\var{rr} & (18925, 60) & The 60 RR intervals of the segment, in seconds \\
\var{af\_frac} & (18925,) & Fraction of the segment's intervals annotated as
AF \\
\var{label\_50} & (18925,) & Main label: 1 = AF, 0 = non-AF \\
\var{label\_excl} & (18925,) & Strict label: 1 / 0 / $-1$ (ambiguous) \\
\var{record} & (18925,) & Patient (record) ID, used for patient-level
splits \\
\var{start\_samp}, \var{end\_samp} & (18925,) & Sample positions of the
first and last beat in the original signal \\
\var{beat\_source} & (18925,) & Which beat file was used (\var{qrs} or
\var{qrsc}) \\
\bottomrule
\end{tabularx}
\end{table}

\begin{table}[ht]
\centering
\caption{Two segments from the same patient (first 8 of 60 intervals shown).}
\label{tab:afdb-sample}
\begin{tabularx}{\textwidth}{@{}cccccL@{}}
\toprule
record & \var{af\_frac} & \var{label\_50} & \var{label\_excl} &
\var{start\_samp} & \var{rr[0:8]} (s) \\
\midrule
04746 & 0.0 & 0 & 0 & 124 &
0.932, 0.904, 0.904, 0.896, 0.920, 0.944, 0.936, 0.928 \\
04746 & 1.0 & 1 & 1 & 2{,}609{,}859 &
0.596, 0.700, 0.480, 0.464, 0.616, 0.676, 0.556, 0.584 \\
\bottomrule
\end{tabularx}
\end{table}

The non-AF segment is regular, with intervals around 0.92\,s. The AF segment
is shorter on average and visibly irregular.

\paragraph{Patient-level folds.} The proportion of AF differs greatly between
patients: two records are 100\% AF, while another is 0.2\% AF. We therefore
split the 24 patients into 5 folds with \var{StratifiedGroupKFold}, grouped
by patient and stratified by label, so that every fold has a similar AF share
(36--46\%). The fold assignment is fixed and shared across all experiments.

%---------------------------------------------------------------------
\subsection{Long-term ECG and wrist PPG with annotated AF episodes (Zenodo)}
%---------------------------------------------------------------------

\paragraph{Source.} Zenodo, version~6 (published 3 November 2026),
\url{https://zenodo.org/records/23921320}, DOI
\href{https://doi.org/10.5281/zenodo.23921320}{10.5281/zenodo.23921320}. It
is open access under the CC BY-NC-SA 4.0 licence, so the data are used for
non-commercial research and are not redistributed in our repository. The
data were collected at Vilnius University and Kaunas University of
Technology.

\paragraph{Content.} 39 patients were monitored for 5--8 days each, 261 days
in total, during their daily life. Each patient wore two devices at the same
time:
\begin{itemize}
  \item a \textbf{reference ECG} (Bittium Faros 180 chest device): single
  lead, 500\,Hz, with a 25\,Hz chest accelerometer;
  \item a \textbf{wrist device}: green-light PPG at 100\,Hz, an ambient-light
  channel, and a 50\,Hz wrist accelerometer.
\end{itemize}
AF episodes were detected automatically on the ECG, then reviewed and
corrected by clinicians. The label is given \emph{for every RR interval}.
Four patients (024, 025, 034, 037) have no PPG data, so 35 patients can be
used for the PPG experiments. The full dataset is 62.7\,GB, so we download
and process a selected subset of patients.

\paragraph{File structure.} MATLAB v7.3 / HDF5, read with \var{h5py}.
Tables~\ref{tab:zen-ecg} and~\ref{tab:zen-ppg} describe the two files of one
patient. In the PPG file, every variable is a list with one entry per
continuous wrist recording (``PPG segment''). A file
\var{subject\_info.xlsx} also gives demographics (sex, age, height, weight),
lifestyle, comorbidities and medication for each patient.

\begin{table}[ht]
\centering
\caption{Variables in \var{XXX\_ECG.mat} (one file per patient).}
\label{tab:zen-ecg}
\begin{tabularx}{\textwidth}{@{}llL@{}}
\toprule
Variable & Shape & Meaning \\
\midrule
\var{ECG} & $(N, 1)$ & ECG signal, 500\,Hz, $\mu$V. Gaps between recording
sessions are filled with NaN \\
\var{Accelerometer\_X/Y/Z} & $(N/20, 1)$ & Chest acceleration, 25\,Hz, mg \\
\var{QRSindex} & $(n_\text{beats}, 1)$ & R-peak sample positions (1-based) \\
\var{rr} & $(n_\text{beats}-1, 1)$ & RR intervals in seconds
($=\mathrm{diff}(\var{QRSindex})/500$) \\
\var{AF\_annotation} & $(n_\text{beats}-1, 1)$ & Label of each RR interval:
1 = AF, 0 = non-AF \\
\shortstack[l]{\var{recording\_startday},\\ \var{recording\_starttime}} & text & Start of the
recording (day number, clock time) \\
\bottomrule
\end{tabularx}
\end{table}

\begin{table}[ht]
\centering
\caption{Variables in \var{XXX\_PPG.mat} (one entry per PPG segment).}
\label{tab:zen-ppg}
\begin{tabularx}{\textwidth}{@{}lL@{}}
\toprule
Variable & Content of each entry \\
\midrule
\var{PPG\_GREEN} & Green-light PPG waveform, 100\,Hz, ADC units \\
\var{PPG\_AMBIENT} & Ambient-light channel, 100\,Hz \\
\var{Accelerometer\_X/Y/Z} & Wrist acceleration, 50\,Hz, g \\
\shortstack[l]{\var{recording\_startday},\\ \var{recording\_starttime}} & Start of the PPG
segment on the ECG clock \\
\var{sync\_method} & How the PPG segment was aligned to the ECG: \var{beat}
(beat-level alignment, about $\pm 0.05$\,s) or \var{coarse} (motion-based
alignment only, about $\pm 1$\,s) \\
\bottomrule
\end{tabularx}
\end{table}

\paragraph{Sample --- patient 006.} The ECG covers 183.3\,h with 864{,}623
beats, 25.3\% of the time in AF (2 AF episodes). The first beat is at
$\var{QRSindex}[0] = 131$, i.e.\ 0.26\,s after the start of the recording.
0.56\% of the RR intervals lie outside $[0.25, 2.5]$\,s, for example long gaps
where the electrodes lost contact. Table~\ref{tab:zen-sample} lists the
patient's PPG segments.

\begin{table}[ht]
\centering
\caption{PPG segments of patient 006.}
\label{tab:zen-sample}
\begin{tabular}{@{}cccccc@{}}
\toprule
PPG segment & Start (from ECG start) & Length & \var{sync\_method} &
ECG beats covered & AF share \\
\midrule
0 & $+18.34$\,h & 2.39\,h & \var{beat} & 9{,}974 & 0\% \\
1 & $+22.08$\,h & 0.80\,h & \var{coarse} & 3{,}274 & 0\% \\
2 & $+45.37$\,h & 2.95\,h & \var{beat} & 17{,}495 & 100\% \\
\bottomrule
\end{tabular}
\end{table}

This sample shows an important property of the dataset: \textbf{PPG covers
only part of the ECG monitoring time} (6.1 of 183.3 hours for patient 006).
The amount of usable labelled PPG is therefore much smaller than the dataset
size suggests.

\paragraph{Transferring ECG labels to PPG.} Labels exist only on the ECG.
Because the PPG start times are given on the ECG clock, each PPG segment is
matched to the ECG beats that fall within its time span, and those beats'
\var{AF\_annotation} values become the labels of the PPG segment. Pulse peaks
detected in \var{PPG\_GREEN} give the PPI series. This series is then cut into
the same 60-interval segments and labelled with the same rule as for AFDB
(\var{label\_50}). We prefer PPG segments with $\var{sync\_method} =
\var{beat}$, since \var{coarse} alignment ($\pm 1$\,s) is too loose for
beat-level labels. The wrist accelerometer is aligned in the same way and is
used in RQ4 to group segments into rest, low-activity and high-activity
conditions.

%---------------------------------------------------------------------
\subsection{Summary of the common representation}
%---------------------------------------------------------------------

\begin{table}[ht]
\centering
\caption{The three domains reduced to one representation.}
\label{tab:common}
\begin{tabularx}{\textwidth}{@{}lLLL@{}}
\toprule
 & AFDB (source) & Zenodo ECG (reference) & Zenodo PPG (target) \\
\midrule
Beat times & Annotated QRS & \var{QRSindex} & Detected pulse peaks \\
Interval series & RR & RR & PPI \\
Label & Cardiologist rhythm annotation & Per-RR \var{AF\_annotation} &
Transferred from simultaneous ECG \\
Unit of analysis & \multicolumn{3}{l}{60-interval segment, AF / non-AF} \\
Split & \multicolumn{3}{l}{Patient-level} \\
\bottomrule
\end{tabularx}
\end{table}

%=====================================================================
\begin{thebibliography}{99}
\small

\bibitem{wolf1991}
P.~A. Wolf, R.~D. Abbott, and W.~B. Kannel.
Atrial fibrillation as an independent risk factor for stroke: the
Framingham Study.
\emph{Stroke}, 22(8):983--988, 1991.

\bibitem{perez2019}
M.~V. Perez et~al.
Large-scale assessment of a smartwatch to identify atrial fibrillation.
\emph{New England Journal of Medicine}, 381(20):1909--1917, 2019.

\bibitem{physionet}
A.~L. Goldberger et~al.
PhysioBank, PhysioToolkit, and PhysioNet: components of a new research
resource for complex physiologic signals.
\emph{Circulation}, 101(23):e215--e220, 2000.

\bibitem{peng1995}
C.-K. Peng, S.~Havlin, H.~E. Stanley, and A.~L. Goldberger.
Quantification of scaling exponents and crossover phenomena in
nonstationary heartbeat time series.
\emph{Chaos}, 5(1):82--87, 1995.

\bibitem{richman2000}
J.~S. Richman and J.~R. Moorman.
Physiological time-series analysis using approximate entropy and sample
entropy.
\emph{American Journal of Physiology -- Heart and Circulatory Physiology},
278(6):H2039--H2049, 2000.

\bibitem{lake2011}
D.~E. Lake and J.~R. Moorman.
Accurate estimation of entropy in very short physiological time series:
the problem of atrial fibrillation detection in implanted ventricular
devices.
\emph{American Journal of Physiology -- Heart and Circulatory Physiology},
300(1):H319--H325, 2011.

\bibitem{bandt2002}
C.~Bandt and B.~Pompe.
Permutation entropy: a natural complexity measure for time series.
\emph{Physical Review Letters}, 88(17):174102, 2002.

\bibitem{costa2002}
M.~Costa, A.~L. Goldberger, and C.-K. Peng.
Multiscale entropy analysis of complex physiologic time series.
\emph{Physical Review Letters}, 89(6):068102, 2002.

\bibitem{takens1981}
F.~Takens.
Detecting strange attractors in turbulence.
In \emph{Dynamical Systems and Turbulence}, Lecture Notes in Mathematics,
vol.~898, pages 366--381. Springer, 1981.

\bibitem{edelsbrunner2010}
H.~Edelsbrunner and J.~Harer.
\emph{Computational Topology: An Introduction}.
American Mathematical Society, 2010.

\bibitem{cohensteiner2007}
D.~Cohen-Steiner, H.~Edelsbrunner, and J.~Harer.
Stability of persistence diagrams.
\emph{Discrete \& Computational Geometry}, 37(1):103--120, 2007.

\bibitem{adams2017}
H.~Adams et~al.
Persistence images: a stable vector representation of persistent homology.
\emph{Journal of Machine Learning Research}, 18(8):1--35, 2017.

\bibitem{bubenik2015}
P.~Bubenik.
Statistical topological data analysis using persistence landscapes.
\emph{Journal of Machine Learning Research}, 16:77--102, 2015.

\bibitem{moody1983}
G.~B. Moody and R.~G. Mark.
A new method for detecting atrial fibrillation using R-R intervals.
\emph{Computers in Cardiology}, 10:227--230, 1983.

\bibitem{bacevicius2026}
J.~Bacevi\v{c}ius et~al.
Long-term electrocardiogram and wrist-based photoplethysmogram recordings
with annotated atrial fibrillation episodes (Version v6) [Data set].
Zenodo, 2026.
\url{https://doi.org/10.5281/zenodo.23921320}.

\end{thebibliography}

\end{document}
